\chapterhead{41}{ORR-2}{Risk Governance}{1}{4}

\lohead{41}{a}{Explain the Basel regulatory expectations for the governance of an
operational risk management framework.}

Basel II regulates operational risk through three pillars. The first sets the
capital floor, the second lets supervisors demand more where Pillar 1 does not
reach, and the third makes the bank tell the market what it is carrying.

\begin{center}
\begin{tikzpicture}[font=\fontsize{8.2}{10.5}\selectfont,
  ptitle/.style={text width=4.5cm,align=center,rounded corners=3pt,
                 inner sep=6pt,text=white,font=\sffamily\bfseries\fontsize{8.6}{10.5}\selectfont},
  pbody/.style={text width=4.5cm,align=left,rounded corners=3pt,
                inner sep=7pt,fill=paleblue,draw=palenavy,line width=0.6pt,
                text=navy}]

\node[fill=navy,text=white,minimum width=16.2cm,minimum height=0.68cm,
      rounded corners=2.5pt,
      font=\sffamily\bfseries\fontsize{8.8}{11}\selectfont] (cap) at (0,0)
      {BASEL II \ \textbullet\ \ THREE PILLARS FOR THE REGULATION OF OPERATIONAL RISK};

\node[ptitle,fill=rust,anchor=north] (t1) at (-5.55,-0.62) {Pillar 1\\[1pt] Regulatory capital};
\node[ptitle,fill=teal,anchor=north] (t2) at (0,-0.62) {Pillar 2\\[1pt] Supervisory review};
\node[ptitle,fill=forest,anchor=north] (t3) at (5.55,-0.62) {Pillar 3\\[1pt] Market discipline};

\node[pbody,anchor=north] (b1) at (t1.south) {
  $\blacktriangleright$~Minimum capital required to meet any unexpected losses from
  credit, market, and operational risks.\\[3pt]
  $\blacktriangleright$~Minimum coverage ratios to manage liquidity risk.\\[3pt]
  $\blacktriangleright$~Basel Committee's \emph{Principles for the Sound Management
  of Operational Risk}.};

\node[pbody,anchor=north] (b2) at (t2.south) {
  $\blacktriangleright$~Extra capital requirements on top of Pillar 1, addressing
  risks not explicitly considered in Pillar 1 (concentration, compliance,
  governance risks).\\[3pt]
  $\blacktriangleright$~Voluntary disclosure and self-assessment, subject to
  regulatory review.};

\node[pbody,anchor=north] (b3) at (t3.south) {
  $\blacktriangleright$~Required quarterly and annual financial and risk
  disclosures by banks.\\[3pt]
  $\blacktriangleright$~Underlying idea is to have greater capital reserves to
  balance greater risks taken.};

\begin{scope}[on background layer]
  \node[fit=(b1)(b2)(b3),inner sep=0pt] (allb) {};
\end{scope}
\node[fill=navy,minimum width=16.2cm,minimum height=0.3cm,rounded corners=2.5pt,
      anchor=north] at ($(allb.south)+(0,-0.12)$) {};
\end{tikzpicture}
\end{center}
\orfigcap{The three pillars of Basel II as they apply to operational risk.}

\begin{orkeybox}
{\sffamily\bfseries\fontsize{8.6}{10}\selectfont\color{navy}Two points that sit
behind all three pillars}\par
\vspace{4pt}
\begin{lettered}
\item Basel II emphasises not only the quantitative aspects of managing
operational risk but also the qualitative factors, such as governance, culture,
and the effectiveness of the operational risk management framework.
\item Regulatory expectations extend beyond mere compliance, requiring that risk
management be integrated into the fabric of banking operations, with a clear
focus on continuous improvement and the adaptability of risk management
practices to evolving challenges.
\end{lettered}
\end{orkeybox}

% ---------------------------------------------------------------------
\lohead{41}{b}{Describe and compare the roles of different committees and the
board of directors in operational risk governance.}

\orsub{Three levels of operational risk committees}

In large banks, operational risk committees are structured across three levels.

\begin{lettered}
\item \textbf{Lowest level:} focuses on specific business areas (like personal
banking or asset management) or regions, providing data for firmwide risk
assessment and escalating significant issues.
\item \textbf{Middle level:} the organisation's risk committee synthesises
information from the lower level to manage overall operational risk, reporting
regularly to executive and board risk committees.
\item \textbf{Top level:} the board risk committee oversees all levels of
operational risk management, advising the board on risk exposures and making key
decisions. Members are experienced in risk management.
\end{lettered}

\orsub{Governance and risk documentation}

\begin{enumerate}
\item[\textcolor{rust}{\sffamily\bfseries 1)}] \textbf{Committee Terms of
Reference (TOR) document.} Defines the committee's mission, objectives,
membership roles, and meeting schedule, and ensures alignment between risk
reporting and decisions. Committees approve policies and procedures related to
ORM.
\item[\textcolor{rust}{\sffamily\bfseries 2)}] \textbf{Policies and procedures.}
Act as internal controls with step-by-step instructions for specific processes,
and serve as training materials for new and existing employees to reduce errors.
\end{enumerate}

\orsub{Board of directors role}

\noindent\begin{minipage}{\textwidth}
\renewcommand{\arraystretch}{1.4}
\begin{tabularx}{\textwidth}{@{}YY@{}}
\rowcolor{navy} \thd{Operational Risk} & \thd{Operational Resilience}\\
{\tabfont a) Sets risk tolerance} & {\tabfont a) Articulates operational resilience goals}\\
\rowcolor{paleblue}
{\tabfont b) Approves and updates the ORMF} & {\tabfont b) Integrates risk tolerance}\\
{\tabfont c) Ensures policy execution by senior management} & {\tabfont c) Monitors senior management}\\
\rowcolor{paleblue}
{\tabfont d) Promotes a culture of risk management} & {\tabfont d) Allocates resources for resilience and ensures training}\\
\end{tabularx}
\end{minipage}
\orfigcap{What the board owns, split across risk and resilience.}

% ---------------------------------------------------------------------
\lohead[Describe the three lines of defense model for operational risk governance
and compare roles and responsibilities for each line of defense.]%
{41}{c}{Describe the ``three lines of defense'' model for operational risk
governance and compare roles and responsibilities for each line of defense.}

\noindent\begin{minipage}{\textwidth}
\renewcommand{\arraystretch}{1.32}
\begin{tabularx}{\textwidth}{@{}YYY@{}}
\rowcolor{navy} \thd{Line 1} & \thd{Line 2} & \thd{Line 3}\\
{\tabfont The business units or risk owners that generate, measure, and manage
risks.} &
{\tabfont Objective review and oversight of Line 1's risk management
activities.} &
{\tabfont Internal audit, separate from risk management, conducting independent
evaluations.}\\
\rowcolor{palenavy}
\multicolumn{3}{@{}l@{}}{\hspace{6pt}\tsub{Roles and responsibilities}}\\
{\tabfont
$\blacktriangleright$~Identify operational risks.\newline
$\blacktriangleright$~Establish controls to manage risks.\newline
$\blacktriangleright$~Evaluate control effectiveness.\newline
$\blacktriangleright$~Provide oversight and reporting on operational risk within
the business line.} &
{\tabfont
$\blacktriangleright$~Develop ORM policies and procedures.\newline
$\blacktriangleright$~Cross-examine Line 1's work, including risk measurement and
reporting.\newline
$\blacktriangleright$~Oversee the bank's monitoring and reporting functions.\newline
$\blacktriangleright$~Provide input on major business decisions and regulatory
compliance.} &
{\tabfont
$\blacktriangleright$~Review controls and adherence to policies and
procedures.\newline
$\blacktriangleright$~Maintain its own list of significant risks.\newline
$\blacktriangleright$~Occasionally share information with Line 2.}\\
\rowcolor{palenavy}
\multicolumn{3}{@{}l@{}}{\hspace{6pt}\tsub{Relationship with other lines}}\\
{\tabfont
$\blacktriangleright$~Inform Line 2 when unable to perform operational risk
duties.\newline
$\blacktriangleright$~Escalate control weaknesses and process weaknesses to
Line 2.} &
{\tabfont
$\blacktriangleright$~Must maintain independence from Line 1.\newline
$\blacktriangleright$~May share information with Line 3 to reduce
redundancies.} &
{\tabfont
$\blacktriangleright$~Must be strictly separate from other departments.\newline
$\blacktriangleright$~Assess reliability before relying on work from other
departments.\newline
$\blacktriangleright$~Provides independent assurance to internal and external
stakeholders.}\\
\end{tabularx}
\end{minipage}
\orfigcap{The three lines of defense, compared across role and relationship.}

% ---------------------------------------------------------------------
\lohead{41}{d}{Explain best practices and regulatory expectations for the
development of a risk appetite for operational risk and for a strong risk
culture.}

\noindent\begin{minipage}{\textwidth}
\renewcommand{\arraystretch}{1.4}
\begin{tabularx}{\textwidth}{@{}YY@{}}
\rowcolor{navy} \thd{Risk appetite --- regulatory expectations} & \thd{Risk appetite --- best practices}\\
{\tabfont a) Define the bank's risk appetite in alignment with strategy} & {\tabfont a) Optimise the risk-return tradeoff}\\
\rowcolor{paleblue}
{\tabfont b) Set clear limits and controls for managing major risks} & {\tabfont b) Establish exposure limits}\\
{\tabfont c) Monitor metrics and adjust based on scenario analysis} & {\tabfont c) Customise appetite for different risks}\\
\rowcolor{paleblue}
{\tabfont d) Ensure consistency between policy and practice} & {\tabfont d) Implement oversight measures}\\
\end{tabularx}
\end{minipage}
\orfigcap{Regulatory expectations set the floor; best practice goes beyond it.}

\orsub{Risk culture}

\begin{lettered}
\item A strong risk culture reduces operational risk, promotes ethical behaviour,
and meets regulatory expectations.
\item Driven by ethical leadership from the board and management, supported by a
clear code of conduct.
\item Aligns compensation with wise risk-taking behaviours and offers
comprehensive risk management training.
\item Encourages a culture of accountability and openness, utilising mechanisms
like whistleblowing for risk control.
\end{lettered}

% =====================================================================
